\documentclass{article}

% Language setting
% Replace `english' with e.g. `spanish' to change the document language
\usepackage[spanish]{babel}

% Set page size and margins
% Replace `letterpaper' with `a4paper' for UK/EU standard size
\usepackage[letterpaper,top=2cm,bottom=2cm,left=3cm,right=3cm,marginparwidth=1.75cm]{geometry}


% Useful packages
\usepackage{float}
\usepackage{amsmath}
\usepackage{graphicx}
\usepackage[hyphens]{url}
\usepackage[colorlinks=true, breaklinks=true]{hyperref}

\linespread{1.9}

\title{PC para Desarrollo de Herramientas para Edición de Videos}

\author{Jesús Romero C.I. 32.590.184  Juan Bermúdez C.I. 31.608.369}

\begin{document}
\maketitle

\section{Tabla de componentes}

\begin{tabular}{|c|p{11cm}|l|c|} \hline
    Componente & Modelo & Precio \\ \hline
    Placa Base & Gigabyte B650 & 156.6\$ \\ \hline
    Procesador & AMD Ryzen 9 9900X & 409.36\$ \\ \hline
    Memoria Ram & DDR5 32GB 6000Mhz & 778.36\$ \\ \hline
    Tarjeta Grafica & Nvidia RTX 5070 12GB & 931.15\$ \\ \hline
    Almacenamiento &  M.2. 2TB 7000MB/s & 482.41\$ \\ \hline
    Case & Lian Li V100 ATX & 75\$ \\ \hline
    Fuente de Poder & Phanteks AMP 850W 80 Plus Gold & 127.6\$ \\ \hline
    Disipador & Nfortec VELA & 31,32\$ \\ \hline
    Pasta Térmica & NOX Hummer Therma Lite 4g & 8.12\$ \\ \hline
    Total &  & 2999.92\$ \\ \hline
\end{tabular}

\section{Definición de cada componente}

% Placa base
\begin{figure}[H]
\centering
\includegraphics[width=1\linewidth]{ARQUI PC/Placa madre.JPG}
\caption{\label{fig:placa}Placa Base Gigabyte B650}
\end{figure}

\url{https://www.pccomponentes.com/placa-base-gigabyte-b650-amd-am5-microatx-b650m-gaming-plus-wifi-ddr5-2xm2-25gbe-wifi-6e-bluetooth-53}

Es la columna vertebral del sistema, encargada de interconectar y permitir la comunicación entre todos los componentes.

\begin{itemize}
    \item \text{\large Tecnologías Clave:} Cuenta con soporte para overclocking tanto en CPU como en RAM. Suele integrar conectividad de red 2.5GbE y la última conectividad USB 3.2 Gen 2x2 (20Gbps) para periféricos de alta velocidad. Además, incluye múltiples ranuras M.2 para almacenamiento NVMe PCIe 5.0/4.0, un códec de audio de alta definición y WiFi 6E integrado.

    \item \text{\large Beneficios para la Edición:} Ofrece un ecosistema preparado para el futuro, con margen para exprimir el hardware mediante overclocking y con soporte para la última generación de almacenamiento ultrarrápido, lo que reduce drásticamente los tiempos de carga de proyectos y archivos de video pesados.
\end{itemize}

% Procesador
\begin{figure}[H]
    \centering
    \includegraphics[width=1\linewidth]{ARQUI PC/Procesador.JPG}
    \caption{\label{fig:procesador}Procesador AMD Ryzen 9 9900X.}
\end{figure}

\url{https://www.pccomponentes.com/procesador-amd-ryzen-9-9900x-44-56ghz}

Es el cerebro del equipo, responsable de ejecutar todas las instrucciones de software, desde el sistema operativo hasta el renderizado más complejo.

\begin{itemize}
    \item \text{\large Tecnologías Clave:} Compatible con PCIe 5.0 para duplicar el ancho de banda de periféricos y memoria DDR5. Su TDP es de 120W, y utiliza el socket AM5, lo que garantiza un soporte de plataforma a largo plazo. Cuenta con una memoria caché total de 76 MB para acelerar el acceso a los datos frecuentes.

    \item \text{\large Beneficios para la Edición:} Con 12 núcleos y 24 hilos, es una auténtica "bestia" para el renderizado y la codificación en paralelo. Su alta frecuencia turbo de hasta 5.6 GHz asegura una capacidad de respuesta excepcional en tareas de un solo núcleo, como la fluidez al trabajar en la línea de tiempo.
\end{itemize}

% RAM
\begin{figure}[H]
\centering
\includegraphics[width=1\linewidth]{ARQUI PC/RAM.JPG}
\caption{\label{fig:RAM}Memoria Ram DDR5 32GB 6000Mhz.}
\end{figure}

\url{https://www.pccomponentes.com/memoria-ram-gskill-flare-x5-f5-6000j2836g16gx2-fx5w-32gb-2x16gb-ddr5-6000mhz-cl28-amd-expo-blanco}

Es la memoria de trabajo ultrarrápida donde la CPU almacena los datos que necesita de forma inmediata para las tareas activas.

\begin{itemize}
    \item \text{\large Tecnologías Clave:} La memoria DDR5 ofrece un ancho de banda significativamente mayor que la generación anterior. La velocidad de 6000 MHz es el "punto dulce" para los procesadores Ryzen 9000, ya que permite una comunicación a máxima velocidad y de forma estable.

    \item \text{\large Beneficios para la Edición:} La combinación de 32 GB de capacidad y alta velocidad es fundamental para manejar proyectos 4K y 8K complejos sin cierres inesperados. Permite tener el software de edición, los plugins, los efectos y las previsualizaciones de alta resolución en la memoria de acceso rápido, ofreciendo una experiencia de trabajo fluida.
\end{itemize}

% Gráfica
\begin{figure}[H]
\centering
\includegraphics[width=1\linewidth]{ARQUI PC/Grafica.JPG}
\caption{\label{fig:Grafica}Tarjeta gráfica Nvidia RTX 5070 12GB.}
\end{figure}

\url{https://www.pccomponentes.com/tarjeta-grafica-inno3d-geforce-rtx-5070-ichill-x3-12gb-gddr7-reflex-2-rtx-ai-dlss4}

Es el coprocesador que acelera drásticamente las tareas de procesamiento paralelo masivo, como los efectos visuales y la codificación de video.

\begin{itemize}
    \item \text{\large Tecnologías Clave:} Cuenta con los núcleos CUDA para renderizado por GPU y los núcleos Tensor para tareas de IA. Su codificador de hardware NVENC de 9ª generación es su función más importante para este proyecto, ya que está optimizado para exportar video en formatos 4:2:2 a una velocidad hasta 6 veces superior a la generación anterior, liberando completamente a la CPU de esa tarea.

    \item \text{\large Beneficios para la Edición:} La principal ventaja es que permite desarrollar y probar el renderizado acelerado por hardware de tus herramientas de edición, logrando que las exportaciones de video en programas como DaVinci Resolve o Premiere Pro sean casi instantáneas. Además, los 12 GB de VRAM son ideales para trabajar con efectos en tiempo real en resoluciones 4K y 8K.
\end{itemize}

% Almacenamiento
\begin{figure}[H]
\centering
\includegraphics[width=1\linewidth]{ARQUI PC/Almacenamiento.JPG}
\caption{\label{fig:almacenamiento}Memoria de almacenamiento M.2. 2TB 7000MB/s.}
\end{figure}

\url{https://www.amazon.com/-/es/Samsung-interna-velocidad-control-MZ-V8P2T0B/dp/B08RK2SR23/ref=sr_1_9?__mk_es_US=ÅMÅŽÕÑ\&crid=3D102G3TA2EG5\&dib=eyJ2IjoiMSJ9.baSemFTGQqqu_tw6h_h8OJyNRWS4ZmGr9LfJ7Kv1bV79mbd1-nhEHEOkR_vxRt9g5VOu3ZumL7HEb26vKQrUqKUWEYRTjiZuRE27gn23qYiQM0gXrtQAxqg1Q4O5fX4bn15SLjhc7b78mCm4nsEJUc9uqAsLjWd1iBiuIpO0Idyi2Dg8cY4s13gWGm5EHXITGV41uQv0zI-hmkUsjrfcdFS_to1XHUZgE7np52LlOXk.TtL0jLS_14hpxsWSwxqnNXftrHL2GplJHhmWqBr5-Gk\&dib_tag=se\&keywords=Tarjeta\%2Bde\%2Bmemoria\%2BHDD\&qid=1779570517\&refinements=p_n_g-1003417077111\%3A24056401011&rnid=24056389011\&sprefix=tarjeta\%2Bde\%2Bmemoria\%2Bhdd\%2Caps\%2C204&sr=8-9\&th=1}

Es la unidad donde se instala el sistema operativo, el software y los proyectos de trabajo.

\begin{itemize}
    \item \text{\large Tecnologías Clave:} Utiliza la interfaz NVMe, diseñada específicamente para SSDs, sobre un bus PCIe 4.0 x4, lo que permite alcanzar velocidades de lectura secuencial de hasta 7,000 MB/s. Es crucial elegir un modelo con memoria TLC y caché DRAM integrada para asegurar un rendimiento sostenido durante las transferencias de archivos grandes.

    \item \text{\large Beneficios para la Edición:} Esta velocidad de lectura es vital para cargar archivos de video RAW de gran tamaño en cuestión de segundos. Los 2 TB ofrecen un amplio espacio para el proyecto activo y una caché de archivos. A diferencia de un SSD sin DRAM, este mantendrá su velocidad incluso al mover cientos de GB, algo común en la edición de video.
\end{itemize}

% CASE
\begin{figure}[H]
\centering
\includegraphics[width=1\linewidth]{ARQUI PC/Case.JPG}
\caption{\label{fig:Case}Case Lian Li V100 ATX.}
\end{figure}

\url{https://www.amazon.com/-/es/Lian-Li-Ventiladores-Preinstalados-Visualización/dp/B0F2T66QC9/ref=sr_1_5?__mk_es_US=ÅMÅŽÕÑ\&dib=eyJ2IjoiMSJ9.xu74ZjLa75AYvTj6qwFHlcn2S6MKgbHbZNK2f_UO_pKpp9I4Hm77tax5-Vd4F3C8ybFEZ3UJZCWLqT-lPH6zkP2Wb-aZvIYjk0uHzcGyFEvxdG3MnXPbaHsWueA-nieBlsNFbeNgr2rVPU4FEn8ye4oRz4JbACAOwC9JZAK3sfwrso5fduWB9vGdJZRqyqsfYfyjSS1_fGelUPeiOs7kbx_h21i2aoac6a8ebpTGAPY.bMUoWAZHq2FUKqiOMFrZbchQfZlyP6e552LEhrqwQjE\&dib_tag=se\&keywords=case\%2Batx\&qid=1779574659\&sr=8-5\&th=1}

Es el chasis que alberga y protege todos los componentes internos, y su diseño es clave para la gestión térmica y la futura ampliabilidad.

\begin{itemize}
    \item \text{\large Tecnologías Clave:} El estándar ATX es el que proporciona el mayor espacio interno. Un buen gabinete no es solo una caja metálica; su diseño incluye un flujo de aire calculado y múltiples puntos para instalar ventiladores y radiadores de refrigeración líquida.

    \item \text{\large Beneficios para la Edición:} Su principal ventaja es el amplio espacio y la flexibilidad. Facilita un montaje limpio y ordenado, mejora la circulación del aire para refrigerar de forma óptima el Ryzen 9 y la RTX 5070, y ofrece espacio de sobra para añadir más unidades de almacenamiento en el futuro. Esto es esencial en un entorno de trabajo donde el almacenamiento de proyectos crece rápidamente.
\end{itemize}

% Fuente de alimentación
\begin{figure}[H]
\centering
\includegraphics[width=1\linewidth]{ARQUI PC/Fuente de poder.JPG}
\caption{\label{fig:Fuente}Fuente de alimentación certificada 850w.}
\end{figure}

\url{https://www.pccomponentes.com/fuente-alimentacion-phanteks-amp-gh-fuente-alimentacion-850w-80-plus-gold-modular-negra}

Es el componente que suministra energía estable y limpia a todos los componentes del sistema.

\begin{itemize}
    \item \text{\large Tecnologías Clave:} La certificación 80 Plus Gold garantiza una eficiencia energética de al menos el 90-92\% bajo cargas típicas, evitando desperdicio en forma de calor. El formato ATX y los 850W de potencia proporcionan suficiente margen de seguridad para los picos de consumo de la CPU y GPU bajo carga máxima.

    \item \text{\large Beneficios para la Edición:} La fiabilidad es su mayor activo. Asegura que todos los componentes reciban la energía que necesitan sin fluctuaciones, especialmente durante largas sesiones de renderizado en las que el sistema estará al 100\% de carga. La alta eficiencia también contribuye a mantener el sistema más fresco y silencioso.
\end{itemize}

% Disipador
\begin{figure}[H]
\centering
\includegraphics[width=1\linewidth]{ARQUI PC/Disipador.JPG}
\caption{\label{fig:disipador}Disipador de CPU 120mm.}
\end{figure}

\url{https://www.pccomponentes.com/nfortec-vela-x-5pipes-ventilador-cpu-120mm}

Es el sistema de refrigeración encargado de extraer el calor de la CPU y expulsarlo fuera de la caja.

\begin{itemize}
    \item \text{\large Tecnologías Clave:} Este tipo de disipador de alto rendimiento utiliza un diseño de doble torre de aletas de aluminio, atravesadas por múltiples tubos de calor de cobre. El ventilador de 120 mm con PWM permite un control dinámico de las revoluciones para equilibrar ruido y refrigeración según la carga de trabajo.

    \item \text{\large Beneficios para la Edición:} Una solución de aire de gama alta es más que suficiente para mantener a raya un procesador de 120W bajo carga sostenida de trabajo. Es una opción fiable, sin riesgos de fugas y con un mantenimiento mínimo.
\end{itemize}

\vspace{3cm}

% Pasta térmica
\begin{figure}[H]
\centering
\includegraphics[width=1\linewidth]{ARQUI PC/Pasta termica.JPG}
\caption{\label{fig:pasta térmica}Pasta térmica.}
\end{figure}

\url{https://www.pccomponentes.com/pasta-termica-nox-hummer-therma-lite-4g-alta-conductividad-y-aplicador}

Es el compuesto conductor que se aplica entre el procesador y el disipador para eliminar cualquier espacio de aire microscópico y maximizar la transferencia de calor.

\begin{itemize}
    \item \text{\large Tecnologías Clave:} Su función es puramente física. Al rellenar las imperfecciones de ambas superficies metálicas, asegura un contacto perfecto. Puede ser de diferentes materiales con distintas conductividades, como silicona, cerámica o metal líquido.

    \item \text{\large Beneficios para la Edición:} Una buena aplicación es el primer paso crítico para una refrigeración eficaz, previniendo el sobrecalentamiento y asegurando que el Ryzen 9 9900X mantenga su máximo rendimiento sin sufrir thermal throttling (bajada de frecuencias por temperatura).
\end{itemize}

\end{document}
